\documentclass[10pt,a4paper]{amsart}
\usepackage[margin=19mm]{geometry}
\usepackage{amsmath,amssymb,mathtools}
\usepackage[scr=rsfs,cal=euler]{mathalfa}
\usepackage{xfrac}
\usepackage[unicode,hidelinks,bookmarks=false]{hyperref}
\newcommand{\ie}{i.e.\ }
\newcommand{\cf}{cf.\ }
\newcommand{\resp}{respectively\ }
\newcommand{\Subsection}{\subsection}
\providecommand{\nobreakdash}{\nobreak\hbox{-}}
\setlength{\emergencystretch}{2em}
\multlinegap0pt
\usepackage{bm}
\usepackage{bbm}
\usepackage{dsfont}
\usepackage{xspace}
\usepackage{cancel}
\usepackage{mathtools}
\usepackage{amsthm}
\makeatletter
\@ifundefined{Theorem}{\newtheorem{Theorem}{Theorem}}{}
\@ifundefined{Proposition}{\newtheorem{Proposition}[Theorem]{Proposition}}{}
\makeatother
\newcommand{\De}{\Delta}
\newcommand{\La}{\Lambda}
\newcommand{\al}{\alpha}
\newcommand{\be}{\beta}
\newcommand{\de}{\delta}
\newcommand{\emp}{\emptyset}
\newcommand{\ep}{\varepsilon}
\newcommand{\la}{\lambda}
\newcommand{\mc}[1]{\mathcal{#1}}
\newcommand{\sm}{\setminus}
\newcommand{\sse}{\subseteq}
\newcommand{\ze}{\zeta}
\renewcommand{\th}{\theta}
\title[Spectral Triples on a non-standard presentation of Effros-Sh]{Spectral Triples on a non-standard presentation of Effros-Shen AF algebras}
\author{Konrad Aguilar, Samantha Brooker, Jack Spielberg}
\date{Source: arXiv:2510.18146v1 (CC BY 4.0)}
\begin{document}
\maketitle

\noindent\emph{Source and licence.} Spectral Triples on a non-standard presentation of Effros-Shen AF algebras. Version arXiv:2510.18146v1 (CC BY 4.0), distributed under the Creative Commons Attribution 4.0 International licence, as recorded in the arXiv licence metadata for this exact version and shown on the abstract page https://arxiv.org/abs/2510.18146v1. Full rights notice: licenses/arxiv-2510.18146-CC-BY-4.0.txt.\par\smallskip

\noindent\textit{Editorial scope.} Counted unit: the Proposition whose source label is \texttt{prop F is refined by (c=beta)} in the article (source0.tex lines 1355--1357, source label \texttt{prop F is refined by (c=beta)}), delivered together with the article's own complete original proof of it (source0.tex lines 1359--1468, one \texttt{proof} environment of 110 lines). No mathematics is written, repaired or re-derived: statement and proof are the article's own text line for line. Required context retained uncounted: none. Every definition, display and label of those lines is retained unchanged. Omitted from this package and not redistributed: 1--1354, 1358--1358, 1469--1739. Nothing inside a retained excerpt is omitted or altered: every retained body is delivered exactly as the article's lines are written, so no omission receipt of any kind accompanies this package. Citations: the retained excerpts carry no citation command at all, so no bibliography entry of the article is transcribed and no reference is redistributed. Reference adaptation: none was needed and none was made; every reference inside the delivered unit resolves inside it, so no unresolved reference or citation is produced. Printed slips kept literal: no printed word, symbol, hypothesis or number of the counted statement or of its proof was altered. Layout and class: the article's own class is \texttt{amsart} with packages including amssymb, amsmath, dsfont, array, longtable, url, enumerate, amsthm, amsfonts, tikz, pgf, graphicx, caption, subcaption; the wrapper is the lane's pinned \texttt{amsart} editorial wrapper at 10pt on A4, which restates the theorem-like environments the retained text needs and adds the article's own macro definitions that the retained text needs (\texttt{\textbackslash De}, \texttt{\textbackslash La}, \texttt{\textbackslash al}, \texttt{\textbackslash be}, \texttt{\textbackslash de}, \texttt{\textbackslash emp}, \texttt{\textbackslash ep}, \texttt{\textbackslash la}, \texttt{\textbackslash mc}, \texttt{\textbackslash sm}, \texttt{\textbackslash sse}, \texttt{\textbackslash th}, \texttt{\textbackslash ze}), with extra wrapper packages (bm, bbm, dsfont, xspace, cancel, mathtools). The delivered numbering is the unit's own. Assets: no retained span contains a figure environment, a graphics-inclusion command, a PDF-page inclusion, a table or a TikZ picture; no image file is copied into this package, the package declares no asset dependency and no image file is redistributed with it. Licence: version arXiv:2510.18146v1 of this article is distributed under the Creative Commons Attribution 4.0 International licence, as recorded in the arXiv licence metadata for this exact version and shown on the abstract page https://arxiv.org/abs/2510.18146v1. A full rights notice is delivered at licenses/arxiv-2510.18146-CC-BY-4.0.txt. Characters: every retained excerpt is pure ASCII, so the delivered engine, XeLaTeX with its default Latin Modern text font, reports no missing character.
\begin{Proposition}\label{prop F is refined by (c=beta)}
    Let $j \geq 1, \ell \geq 0, m = m(j,\ell), j' \geq j, \ell' \ge m,$ and $m' = m(j', \ell')$. Let $\De = (\de, \ep) =(\mu\al^r, \nu\be^s) $ with $|\de| = |\ep| = p_j + \ell$, and $\De' = (\de', \ep') = (\nu\be^{r'}, \nu' \al^{s'}) \in S_{j',\ell'}$, where $(\mu, \nu) \in S_j$ and $(\nu, \nu') \in S_{j'}$. Let $n \geq 0$ and fix $E \in Q_n^{m'}$.  Then the collection $\mc{F} = \{\al^r F : F \in E_{n+1}^{m'} \}$ is refined by $Q_{n+m}^{m'-r}$.
\end{Proposition}

\begin{proof}
    Let $m'' = m'-r$ and $n' = n+r$. For any $G \sse Z_{m'}$, let $G' = \al^r G$.
    
    If $E \in Q^{m'}_n \sm P_{m', n}$, so $E = \th R$ for some $\th \in \Phi_n^{m'}$ with $|\th| \ge 1$, and some $R \in P_{m'+|\th|, n-|\th|}$, then $E \cap Q_{n+1}^{m'} \sse Z(\th) \cap Q_{n+1}^{m'} = \th Q_{n+1 - |\th|}^{m'+|\th|}$. And since $\th \in \Phi_n^{m'}$, $\al^r \th \in \Phi_{n'}^{m''}$. Hence, $\mc{F} = \al^r (E^{m'}_{n+1}) \sse \al^r \th Q_{n+1 -|\th|}^{m'+|\th|} = Z(\al^r\th) \cap Q_{n'+1}^{m''}$. In other words, $\mc{F}$ is refined by $Q_{n'+1}^{m''}$.

    Otherwise, $E \in P_{m',n}$. This breaks down into several cases, each of which may also involve several cases...

    \textbf{Case 1:} $E = Z_{m'}(\al^{n+1}\be^J) \sm Z_{m'}(\al^{n+1}\be^{J+1}) \in P_{m',n}^{(1)}$ for some $J \leq n$. Then
    \begin{align*}
        E \cap (Q^{m'}_{n+1} \sm P_{m', n+1}) &= \emp\\
        E \cap P^{(1)}_{m',n+1} &= \{ Z_{m'}(\al^{n+2}\be^J) \sm Z_{m'}(\al^{n+2} \be^{J+1}) \}\\
        E \cap P^{(2)}_{m',n+1} &= \emp\\
        E \cap P^{(3)}_{m',n+1} &= \{ Z_{m'}(\al^{n+1}\be^J \la) : \la \in \La_2^1\} \\
        E \cap P^{(4)}_{m',n+1} &= \emp
    \end{align*}
    \begin{enumerate}
        \item If $F \in P^{(1)}_{m',n+1}$, then $F' = Z_{m''}(\al^{n'+2}\be^{J}) \sm Z_{m''}(\al^{n'+2}\be^{J+1}) \in P^{(1)}_{m'', n' + 1}$.

    \item If $F \in P^{(3)}_{m',n+1}$, then $F' = Z_{m''}(\al^{n'+1}\be^{J} \la) \in P^{(3)}_{m'', n' + 1}$.
    \end{enumerate}
    Thus $\mc{F}$ is refined by $Q^{m''}_{n'+1}$, which is refined by $Q^{m''}_{n+m}$.

    \textbf{Case 2:} $E = Z_{m'}(\al^{I}\be^{n+1}) \sm Z_{m'}(\al^{I+1}\be^{n+1}) \in P_{m',n}^{(2)}$ for some $I \leq n$. Then
    \begin{align*}
        E \cap (Q^{m'}_{n+1} \sm P_{m', n+1}) &= \emp\\
        E \cap P^{(1)}_{m',n+1} &= \emp \\
        E \cap P^{(2)}_{m',n+1} &= \{ Z_{m'}(\al^{I}\be^{n+2}) \sm Z_{m'}(\al^{I+1} \be^{n+2}) \}\\
        E \cap P^{(3)}_{m',n+1} &= \{ Z_{m'}(\al^{I}\be^{n+1} \la) : \la \in \La_2^1\} \\
        E \cap P^{(4)}_{m',n+1} &= \emp
    \end{align*}
    \begin{enumerate}
        \item If $F \in P^{(2)}_{m',n+1}$, then $F' = Z_{m''}(\al^{I+r}\be^{n+2})\sm Z_{m''}(\al^{I+r+1}\be^{n+2})$.
    \begin{enumerate}
        \item If $I+r \leq n+1$, then $F' \in P^{(2)}_{m'', n+1}$.
        \item Else, $n+1 < I+r$, and \begin{align*}
            F' = \left(\bigcup_{h=n+2}^{I+r} \bigcup_{\la \in \La_2^1} Z_{m''}(\al^{I+r}\be^h \la) \right) \cup Z_{m''}(\al^{I+r} \be^{I+r+1}) \sm Z_{m''}(\al^{I+r+1}\be^{I+r+1}).
        \end{align*}
        For $n+2 \leq h \leq I+r$ and $\la \in \La_2^1$, $Z_{m''}(\al^{I+r}\be^h \la) \in P^{(3)}_{m'', I+r}$. And $Z_{m''}(\al^{I+r} \be^{I+r+1}) \sm Z_{m''}(\al^{I+r+1}\be^{I+r+1}) \in P^{(2)}_{m'', I+r}$. Hence, $F'$ is refined by $Q_{I+r}^{m''}$. Since $I \leq n$ and $r \leq m$, $I+r \leq n+m$, so $F'$ is refined by $Q^{m''}_{n+m}$.
    \end{enumerate}

    \item If $F \in P^{(3)}_{m',n+1}$, then $F' = Z_{m''}(\al^{I+r}\be^{n+1}\la)$.
    \begin{enumerate}
        \item If $I+r \leq n+1$, $\al^r F \in P^{(3)}_{m'', n+1}$, which is refined by $Q_{n+m}^{m''}$.
        
        \item Otherwise, $F' \in P_{m'', I+r}^{(3)}$, and $I+r \leq n+m$, so $F'$ is refined by $Q_{n+m}^{m''}$.
    \end{enumerate}
    \end{enumerate}
    Thus $\mc{F}$ is refined by $Q^{m''}_{n+m}$.

    \textbf{Case 3:} If $E = Z_{m'}(\al^I \be^J \la) \in P^{(3)}_{m',n}$ for some $I, J \leq n \leq I+J$ and $\la \in \La_2^1$, let $\ze = \al^I \be^J \la$. Then
    \begin{align*}
        E \cap (Q^{m'}_{n+1} \sm P_{m', n+1}) &= \begin{cases}
            \emp &: I+J \geq n+1\\
            \ze P_{m'+n+1, 0} &: I+J=n
        \end{cases}\\
        E \cap P^{(1)}_{m',n+1} &= \emp\\
        E \cap P^{(2)}_{m',n+1} &= \emp\\
        E \cap P^{(4)}_{m',n+1} &= \emp \\
        E \cap P^{(3)}_{m',n+1} &= \begin{cases}
            \{E\} &: I+J \geq n+1\\ \emp &: I+J = n.
        \end{cases}\\
    \end{align*}
    \begin{enumerate}
        \item First suppose $I+J \ge n+1$. Then if $F \in E_{n+1}^{m'}$, $F = E = Z_{m'}(\al^I \be^J \la)$. So, $F' = Z_{m''}(\al^{I+r}\be^J \la)$. If $I+r \le n+1$, then $F' \in P^{(3)}_{m'', n+1}$, refined by $Q_{n+1}^{m''}$. If $n+1 < I+r$, then $F' \in P^{(3)}_{m'',I+r}$, and so is refined by $Q^{m''}_{I+r}$. Since $I+r \leq n+m$, $F'$ is refined by $Q_{n+m}^{m''}$.

        \item Suppose $I+J = n$. Then if $F \in E_{n+1}^{m'}$, $F \in \ze P_{m'+n+1, 0}$. Since $|\ze| = n + 1$, $\ze \in \Phi_{n+1}^{m'}$, and $\al^r \ze \in \Phi_{n'+1}^{m''}$. Hence, $\al^r \ze P_{m'+n+1,0} \sse Q_{n'+1}^{m''}$, and so $\mc{F}$ is refined by $Q_{n'+1}^{m''}$. Since $n'+1 \leq n+m$, $F'$ is refined by $Q_{n+m}^{m''}$.
    \end{enumerate}
    Therefore, $\mc{F}$ is refined by $Q_{n+m}^{m''}$.

    \textbf{Case 4:} $E = Z_{m'}(\al^{n+1}\be^{n+1}) \in P_{m',n}^{(4)}$. Then
    \begin{align*}
        E \cap (Q^{m'}_{n+1} \sm P_{m', n+1}) &= \emp\\
        E \cap P^{(1)}_{m',n+1} &= \{Z_{m'}(\al^{n+2}\be^{n+1}) \sm Z_{m'}(\al^{n+2}\be^{n+2})\} \\
        E \cap P^{(2)}_{m',n+1} &= \{ Z_{m'}(\al^{n+1}\be^{n+2}) \sm Z_{m'}(\al^{n+2} \be^{n+2}) \}\\
        E \cap P^{(3)}_{m',n+1} &= \{ Z_{m'}(\al^{n+1}\be^{n+1} \la) : \la \in \La_2^1\} \\
        E \cap P^{(4)}_{m',n+1} &= \{Z_{m'}(\al^{n+2}\be^{n+2})\}
    \end{align*}
    \begin{enumerate}
        \item If $F \in P^{(1)}_{m',n+1}$: \begin{align*}
            F' &= Z_{m''}(\al^{n'+2} \be^{n+1})\sm Z_{m''}(\al^{n'+2} \be^{n+2})\\
            &\in P^{(1)}_{m'', n'+1},
        \end{align*}
        so $F'$ is refined by $Q_{n'+1}^{m''}$.

        \item If $F \in P^{(2)}_{m',n+1}$:
        \begin{align*}
            F' &= Z_{m''}(\al^{n'+1}\be^{n+2}) \sm Z_{m''}(\al^{n'+2} \be^{n+2})\\
            &= \left(\bigcup_{h=n+2}^{n'+1} \bigcup_{\la \in \La_2^1} Z_{m''}(\al^{n'+1} \be^h \la) \right) \cup \left( Z_{m''}(\al^{n'+1} \be^{n'+2}) \sm Z_{m''}(\al^{n'+2} \be^{n'+2})\right).
        \end{align*}
        For $n+2 \le h \le n'+1$ and $\la \in \La_2^1$, $Z_{m''}(\al^{n'+1} \be^h \la) \in P^{(3)}_{m'', n'+1}$. And $Z_{m''}(\al^{n'+1} \be^{n'+2}) \sm Z_{m''}(\al^{n'+2} \be^{n'+2}) \in P^{(2)}_{m'', n'+1}$.

        Hence $F'$ is refined by $Q^{m''}_{n'+1}$.

        \item If $F \in P^{(3)}_{m',n+1}$: \begin{align*}
            F' &= Z_{m''}(\al^{n'+1}\be^{n+1}\la)\\
            &\in P^{(3)}_{m'', n'+1}\\
            &\sse Q_{n'+1}^{m''}.
        \end{align*}

        \item If $F \in P^{(4)}_{m',n+1}$: 
        \begin{align*}
            F' &= Z_{m''}(\al^{n'+2} \be^{n+2} ) \\
            &= \left( \bigcup_{h=n+2}^{n'+1} Z_{m''}(\al^{n'+2} \be^h) \sm Z_{m''}(\al^{n'+2}\be^{h+1}) \right) \cup Z_{m''}(\al^{n'+2} \be^{n'+2}).
        \end{align*}
        For $n+2 \le h \le n'+1$, $Z_{m''}(\al^{n'+2} \be^h) \sm Z_{m''}(\al^{n'+2}\be^{h+1}) \in P^{(1)}_{m'', n'+1}$. And $Z_{m''}(\al^{n'+2} \be^{n'+2}) \in P^{(4)}_{m'', n'+1}$.
        
        Hence $F'$ is refined by $Q_{n'+1}^{m''}$.
    \end{enumerate}
    Therefore, in all cases, $\mc{F}$ is refined by $Q_{n+m}^{m''}$.
\end{proof}

\end{document}
